\chapter{The Code Page}\label{cha:codepage}\markboth{The Code Page}{The Code Page}

Every character the K4510 shows is one byte, and the code page says which
character each of the 256 bytes is. The machine's is IBM's code page 437
--- the PC's, the one with the box drawing --- as every PC program and
every bulletin board expects it. That is what it starts in, and what this
book's pages on the machine are written in.

There is a second page, the \emph{K4510 page}, for when it is wanted. It
is CP437 with 26 of its places given to letters Western Europe needs and
CP437 never had: the accented capitals of French, Italian and Spanish,
\texttt{œ} and \texttt{Œ}, the euro, the section and pilcrow signs,
German's quotes, and \texttt{ø}. The places they took were CP437's Greek
and mathematical symbols and a few marks nobody used. The box drawing, the
shades and blocks, and every accented letter CP437 already had --- the
lower-case letters French, Spanish and German are written with --- are in
the same places in both.

\verb|CODEPAGE K4510| chooses it and \verb|CODEPAGE 437| goes back;
\verb|CODEPAGE| alone says which is in use. F12 $\rightarrow$ Terminal
$\rightarrow$ \emph{Code page} is the same choice, and the machine
remembers it. The screen changes at once: the page's font takes the
place of the other, and the characters already on the screen are drawn
again from it.

Like the register reference, this appendix is not written by hand:
\pth{doc/guide/mkcodepage.py} makes it from \pth{core/codepage.h}, which
holds both pages. JIM turning UTF-8 into bytes and back, the keyboard, the
fonts and \texttt{k4510-remote} all read that file, so none of them can
disagree with this page, or with each other.

\section{The K4510 page}
\input{generated/codepage}

\section{The languages each writes}
CP437 has every lower-case letter of French, Spanish, German, Italian and
the Nordic languages, and \texttt{É Ç Ñ Ä Ö Ü Å Æ} among the capitals ---
enough for most written French, which leaves the other capitals
unaccented. With the K4510 page these are complete in both cases: French,
Italian, Spanish, German, Danish, Norwegian, Swedish, Finnish, Dutch,
Catalan, Galician, Basque, Irish, Scottish Gaelic, Afrikaans and
Albanian. Portuguese is four letters short in either (\texttt{ã õ Ã Õ}).

Text arriving as UTF-8 --- from \texttt{!}, from TELNET, over ssh --- is
turned into the page in use on the way in; a character it has no place for
becomes its nearest, and the curly quotes their plain forms.

\section{The first 32, and \$7F}
Bytes \texttt{\$01}--\texttt{\$1F} and \texttt{\$7F} draw CP437's little
pictures in both pages: the faces, the card suits, the arrows, the notes.
As text they are also the control codes --- \texttt{\$07} rings,
\texttt{\$1B} is Escape, \texttt{\$18} to \texttt{\$1A} move the cursor ---
and the console obeys them rather than drawing them. So a program that
wants a picture on the screen puts the byte into screen memory itself, as
the games do.

\section{Typing them}
Every character a keyboard layout has a key for reaches the machine, if
the page in use has it. The dead keys compose the rest: \texttt{`} then
\texttt{e} is \texttt{è} in either page; \texttt{`} then \texttt{A} is
\texttt{À} in the K4510 page, and a plain \texttt{A} in CP437, which has
no place for it. FONTED redraws any of the 256, in both sizes of the font,
and shows the Unicode value of the character it is editing.

\section{A BBS}
TELNET draws a CP437 session --- a bulletin board, an old system, a far
end that never negotiates --- with strict CP437 whatever page the machine
is in, so ANSI art is always exactly as it was drawn. A Unix host, which
takes the XTERM-COLOR terminal type and talks UTF-8, is drawn in the page
in use, and the machine's font comes back when TELNET hangs up.
